% !TEX program = xelatex
% Source: F:\Obsidian Storage\数学\泛函分析\泛函分析讲义上册习题册.md
\documentclass[UTF8,a4paper,11pt,fontset=fandol]{ctexart}
\usepackage[margin=2.3cm,headheight=15pt]{geometry}
\usepackage{amsmath,amssymb}
\usepackage{fancyhdr}
\usepackage[hidelinks]{hyperref}
\hypersetup{pdftitle={第一, 二周作业},pdfauthor={黄子昱}}
\pagestyle{fancy}
\fancyhf{}
\fancyhead[L]{泛函分析\quad 第一、二周作业}
\fancyhead[R]{黄子昱\quad 24343058}
\fancyfoot[C]{\thepage}
\setlength{\parindent}{2em}
\setlength{\parskip}{0.3em}
\linespread{1.15}
\allowdisplaybreaks[2]
\raggedbottom
\setlength{\emergencystretch}{2em}
\ctexset{section={format=\large\bfseries,beforeskip=1.3ex,afterskip=.7ex},subsection={format=\normalsize\bfseries,beforeskip=1.1ex,afterskip=.5ex}}
\newcommand{\exercise}[1]{\par\ifdim\dimexpr\pagegoal-\pagetotal\relax<7\baselineskip\newpage\fi\subsection*{习题 #1}}
\newcommand{\answer}{\par\smallskip\noindent\textbf{答案}\quad}
\newcommand{\qedherecustom}{\unskip\hfill\mbox{\normalfont Q.E.D.}\par}
\begin{document}
\begin{center}
{\LARGE\bfseries 第一、二周作业}\par\medskip
{\large 泛函分析}\par\medskip
姓名：黄子昱\qquad 学号：24343058
\end{center}
\medskip

\section*{1.1 压缩映象原理}

\exercise{1.1.1}

证明完备空间的闭子集是一个完备的子空间，而任一度量空间中的完备子空间必是闭子集。

\answer

（1）设 $(\mathcal X,\rho)$ 是一个完备的度量空间，取闭集 $M\subset\mathcal X$。要证 $M$ 是一个完备子空间，即证：对任意序列 $\{x_n\}\subset M$，
\[
\rho(x_m,x_n)\longrightarrow0\quad(m,n\to\infty)
\quad\Longrightarrow\quad
\exists x\in M,\quad x_n\longrightarrow x.
\]
由于 $\mathcal X$ 完备，存在 $x\in\mathcal X$，使 $x_n\to x$。而 $M$ 是闭集，包含其中所有收敛序列在 $\mathcal X$ 中的极限，所以 $x\in M$。因此 $M$ 是一个完备子空间。

（2）设 $M$ 是任意度量空间 $(\mathcal X,\rho)$ 的一个完备子空间。要证 $M$ 是闭集，即证
\[
x_n\in M,\quad x_n\longrightarrow x\in\mathcal X
\quad\Longrightarrow\quad x\in M.
\]
由三角不等式，
\[
\rho(x_m,x_n)\leq\rho(x_m,x)+\rho(x,x_n)\longrightarrow0,
\]
所以收敛序列 $\{x_n\}$ 是基本列。由 $M$ 完备，存在 $y\in M$，使 $x_n\to y$。再由
\[
\rho(x,y)\leq\rho(x,x_n)+\rho(x_n,y)\longrightarrow0,
\]
得 $x=y\in M$。因此 $M$ 是闭集。空集的情形结论也成立。

综上所述，\qedherecustom

\exercise{1.1.2}

设 $f$ 定义在 $[a,b]$ 上, 是二阶连续可微的实值函数, $\hat{x}\in(a,b)$, $f(\hat{x})=0$, 且 $f'(\hat{x})\ne0$. 证明存在 $\hat{x}$ 的邻域 $U(\hat{x})$, 使得任意 $x_0\in U(\hat{x})$ 生成的迭代

\[
x_{n+1}=x_n-\frac{f(x_n)}{f'(x_n)}
\]

均收敛, 且 $\lim_{n\to\infty}x_n=\hat{x}$.

\answer

定义映射
\[
g(x)=x-\frac{f(x)}{f'(x)}.
\]
由 $f'(\hat x)\ne0$ 及 $f'$ 连续，可先取 $\hat x$ 的一个邻域，使其中 $f'(x)\ne0$，从而 $g$ 有定义且连续可微。计算得
\[
g'(x)=1-\frac{(f'(x))^2-f(x)f''(x)}{(f'(x))^2}
=\frac{f(x)f''(x)}{(f'(x))^2}.
\]
因为 $f(\hat x)=0$，所以 $g(\hat x)=\hat x$ 且 $g'(\hat x)=0$。由 $g'$ 连续，可取 $r>0$ 和 $0<\alpha<1$，使
\[
I=[\hat x-r,\hat x+r]\subset(a,b),\qquad
|g'(x)|\leq\alpha\quad(x\in I).
\]
由中值定理，对任意 $x,y\in I$，有
\[
|g(x)-g(y)|\leq\alpha|x-y|.
\]
特别地，$|g(x)-\hat x|\leq\alpha|x-\hat x|\leq\alpha r<r$，所以 $g(I)\subset I$。闭区间 $I$ 完备，因此由压缩映射原理，$g$ 在 $I$ 中有唯一不动点，且由任意 $x_0\in I$ 出发的迭代均收敛到它。由于 $\hat x$ 已是不动点，取 $U(\hat x)=(\hat x-r,\hat x+r)$，就有
\[
|x_n-\hat x|\leq\alpha^n|x_0-\hat x|\longrightarrow0.
\]
因此 $x_n\to\hat x$。\qedherecustom

\exercise{1.1.3}

设 $(\mathcal X,\rho)$ 是度量空间, $T:\mathcal X\to\mathcal X$ 满足当 $x\ne y$ 时

\[
\rho(Tx,Ty)<\rho(x,y),
\]

且 $T$ 有不动点. 证明该不动点唯一.

\answer

反证法。若有两个不同的不动点，记为 $x_1\ne x_2$，则 $Tx_1=x_1$，$Tx_2=x_2$。由题设严格缩短距离的条件，有
\[
\rho(x_1,x_2)=\rho(Tx_1,Tx_2)<\rho(x_1,x_2),
\]
矛盾。因此不动点唯一。\qedherecustom

\exercise{1.1.4}

证明度量空间上的压缩映象是连续的.

\answer

设 $T$ 的压缩常数为 $0\leq\alpha<1$，即
\[
\rho(Tx,Ty)\leq\alpha\rho(x,y).
\]
固定任意 $x_0$。对任意 $\varepsilon>0$，取 $\delta=\varepsilon$，则当 $\rho(x,x_0)<\delta$ 时，
\[
\rho(Tx,Tx_0)\leq\alpha\rho(x,x_0)
\leq\rho(x,x_0)<\varepsilon.
\]
因此 $T$ 在任意 $x_0$ 处连续，即 $T$ 连续。\qedherecustom

\exercise{1.1.5}

设 $T$ 是压缩映象. 证明对 $n\in\mathbb N$, $T^n$ 也是压缩映象, 并说明逆命题不一定成立.

\answer

先证原命题。设 $T$ 的压缩常数为 $0\leq\alpha<1$，则反复使用压缩不等式，得到对每个正整数 $n$，
\[
\rho(T^nx,T^ny)\leq\alpha\rho(T^{n-1}x,T^{n-1}y)
\leq\cdots\leq\alpha^n\rho(x,y).
\]
由于 $0\leq\alpha^n<1$，所以 $T^n$ 也是压缩映射。

现在说明逆命题不一定成立。设 $\mathcal X=\mathbb R^2$，取通常的欧氏距离，定义 $T(x,y)=(2y,0)$。对任意 $(x,y)$，
\[
T^2(x,y)=T(2y,0)=(0,0),
\]
所以 $T^2$ 是零映射，为压缩映射。而取 $u=(0,1)$、$v=(0,0)$，有
\[
\|Tu-Tv\|_2=2>1=\|u-v\|_2,
\]
所以 $T$ 不为压缩映射。因此逆命题不成立。\qedherecustom

\exercise{1.1.6}

设 $M$ 是 $(\mathbb R^n,\rho)$ 中的有界闭集, $T:M\to M$, 且 $x\ne y$ 时有 $\rho(Tx,Ty)<\rho(x,y)$. 证明 $T$ 有唯一不动点.

\answer

设 $M\ne\varnothing$，并取 $\mathbb R^n$ 上通常的欧氏距离。由 1.1.1 知 $M$ 完备；又因为 $M$ 有界且闭，所以 $M$ 是紧集。

题设仅给出不同点之间的距离严格缩短，未给出统一的小于 $1$ 的压缩常数，因此还需利用紧性证明不动点存在。由题设，对任意 $x,y\in M$，都有
\[
\rho(Tx,Ty)\leq\rho(x,y),
\]
故 $T$ 连续。令 $h(x)=\rho(x,Tx)$，则
\[
|h(x)-h(y)|\leq\rho(x,y)+\rho(Tx,Ty)\leq2\rho(x,y),
\]
所以 $h$ 连续，并在非空紧集 $M$ 上取到最小值。设最小值点为 $x_0$。

若 $Tx_0\ne x_0$，则由题设，
\[
h(Tx_0)=\rho(Tx_0,T^2x_0)<\rho(x_0,Tx_0)=h(x_0).
\]
但 $Tx_0\in M$，这与 $h(x_0)$ 为最小值矛盾。因此 $Tx_0=x_0$，不动点存在。再由 1.1.3 知不动点唯一。\qedherecustom

\exercise{1.1.7}

设 $y\in C[0,1]$, $|\lambda|<1$. 证明积分方程

\[
x(t)-\lambda\int_0^1e^{t-s}x(s)\,ds=y(t)
\]

在 $C[0,1]$ 中有唯一解.

\answer

原方程等价于
\[
e^{-t}x(t)-\lambda\int_0^1e^{-s}x(s)\,ds=e^{-t}y(t).
\]
作变换，令 $z(t)=e^{-t}x(t)$、$\zeta(t)=e^{-t}y(t)$，定义映射
\[
(Tz)(t)=\zeta(t)+\lambda\int_0^1z(s)\,ds.
\]
由于 $\zeta$ 连续，而积分项关于 $t$ 为常数，所以 $T$ 将 $C[0,1]$ 映入自身。在上确界范数下，对任意 $z_1,z_2\in C[0,1]$，有
\[
\begin{aligned}
\|Tz_1-Tz_2\|_\infty
&=|\lambda|\left|\int_0^1(z_1(s)-z_2(s))\,ds\right|\\
&\leq|\lambda|\int_0^1|z_1(s)-z_2(s)|\,ds
\leq|\lambda|\|z_1-z_2\|_\infty.
\end{aligned}
\]
因为 $|\lambda|<1$ 且 $C[0,1]$ 完备，由压缩映射原理，$T$ 有唯一不动点 $z$。变换 $x(t)=e^tz(t)$ 在 $C[0,1]$ 上可逆，因此原积分方程有唯一解。\qedherecustom

\section*{1.2 完备度量空间}

\exercise{1.2.1}

设 $S$ 是所有实数或复数列 $x=(\xi_1,\xi_2,\ldots,\xi_n,\ldots)$ 的集合. 对 $x=(\xi_k)$ 和 $y=(\eta_k)$ 定义

\[
\rho(x,y)=\sum_{k=1}^{\infty}\frac{1}{2^k}\frac{|\xi_k-\eta_k|}{1+|\xi_k-\eta_k|}.
\]

证明 $(S,\rho)$ 是完备度量空间.

\answer

先证 $\rho$ 是一个距离。记底域为 $\mathbb K=\mathbb R$ 或 $\mathbb C$, 并令
\[
\varphi(t)=\frac{t}{1+t},\qquad t\geq0.
\]
由于 $\varphi'(t)=1/(1+t)^2>0$, 函数 $\varphi$ 在 $[0,\infty)$ 上严格递增。因为 $0\leq\varphi(t)<1$, 所以
\[
0\leq\frac1{2^k}\varphi\left(|\xi_k-\eta_k|\right)\leq\frac1{2^k},
\qquad
\sum_{k=1}^{\infty}\frac1{2^k}=1,
\]
故定义中的级数总是收敛。正定性和对称性显然成立。若 $\rho(x,y)=0$, 由于各项均非负, 每一项都必须为零, 从而 $\xi_k=\eta_k$ 对所有 $k$ 成立。

再证三角不等式。对 $x=(\xi_k)$, $y=(\eta_k)$ 和 $z=(\zeta_k)$, 固定 $k$, 令
\[
a_k=\xi_k-\eta_k,\qquad b_k=\eta_k-\zeta_k.
\]
由 $\varphi(t)=1-1/(1+t)$, 对任意 $s,t\geq0$ 有
\[
\begin{aligned}
\varphi(s+t)
&=1-\frac1{1+s+t}\\
&\leq1-\frac1{(1+s)(1+t)}\\
&=1-(1-\varphi(s))(1-\varphi(t))\\
&=\varphi(s)+\varphi(t)-\varphi(s)\varphi(t)\\
&\leq\varphi(s)+\varphi(t).
\end{aligned}
\]
其中第二行使用了 $(1+s)(1+t)=1+s+t+st\geq1+s+t$, 最后一行使用了 $\varphi(s),\varphi(t)\geq0$。另一方面, 因为 $\xi_k-\zeta_k=a_k+b_k$, 由标量距离的三角不等式,
\[
|\xi_k-\zeta_k|=|a_k+b_k|
\leq|a_k|+|b_k|.
\]
于是
\[
\begin{aligned}
\varphi(|\xi_k-\zeta_k|)
&=\varphi(|a_k+b_k|)\\
&\leq\varphi(|a_k|+|b_k|)\\
&\leq\varphi(|a_k|)+\varphi(|b_k|)\\
&=\varphi(|\xi_k-\eta_k|)+\varphi(|\eta_k-\zeta_k|).
\end{aligned}
\]
乘以 $2^{-k}$ 并对 $k$ 求和, 得
\[
\rho(x,z)\leq\rho(x,y)+\rho(y,z).
\]
因此 $\rho$ 是 $S$ 上的度量。

再证完备性。设 $\{x^{(m)}\}$ 是 $S$ 中的基本列, 对每个 $m$ 按坐标写成
\[
x^{(m)}=(\xi^{(m)}_1,\xi^{(m)}_2,\ldots)\in S.
\]

固定 $k$ 和 $\delta>0$。因为 $\varphi(\delta)>0$, 由基本列的定义, 存在 $N$, 使得当 $m,n\geq N$ 时
\[
\rho(x^{(m)},x^{(n)})<\frac1{2^k}\varphi(\delta).
\]
于是
\[
\varphi\left(|\xi^{(m)}_k-\xi^{(n)}_k|\right)
\leq2^k\rho(x^{(m)},x^{(n)})
<\varphi(\delta).
\]
由 $\varphi$ 严格递增, 得
\[
|\xi^{(m)}_k-\xi^{(n)}_k|<\delta.
\]
所以对每个固定的 $k$, $\{\xi^{(m)}_k\}_m$ 是 $\mathbb K$ 中的标量基本列。

由于 $\mathbb R$ 和 $\mathbb C$ 完备, 对每个 $k$ 都存在标量 $\xi_k^*$, 使
\[
\xi^{(m)}_k\longrightarrow\xi_k^*\qquad(m\to\infty).
\]
构造
\[
x^*=(\xi_1^*,\xi_2^*,\ldots)\in S.
\]
给定 $\varepsilon>0$, 取 $K$ 使
\[
\sum_{k=K+1}^{\infty}\frac1{2^k}<\frac\varepsilon2.
\]
由有限多个坐标的逐一收敛,
\[
\sum_{k=1}^{K}\frac1{2^k}
\varphi\left(|\xi^{(m)}_k-\xi_k^*|\right)\longrightarrow0.
\]
故当 $m$ 足够大时, 上式小于 $\varepsilon/2$。于是
\[
\begin{aligned}
\rho(x^{(m)},x^*)
&\leq\sum_{k=1}^{K}\frac1{2^k}
\varphi\left(|\xi^{(m)}_k-\xi_k^*|\right)
 +\sum_{k=K+1}^{\infty}\frac1{2^k}<\varepsilon.
\end{aligned}
\]
所以 $x^{(m)}\to x^*$。

这证明了 $S$ 的完备性，所以 $(S,\rho)$ 完备。\qedherecustom

\exercise{1.2.2}

证明度量空间中的基本序列收敛, 当且仅当它含有一个收敛子序列.

\answer

必要性显然。若基本列 $\{x_n\}$ 收敛于 $x$, 则它的任意子列都收敛于 $x$, 因而含有收敛子列。

再证充分性。设 $\{x_n\}$ 是 $(\mathcal X,\rho)$ 中的基本列, 且存在子列 $\{x_{n_j}\}$ 收敛于 $x\in\mathcal X$。给定 $\varepsilon>0$, 存在 $N_1$, 使得当 $m,n\geq N_1$ 时
\[
\rho(x_m,x_n)<\frac\varepsilon2.
\]
又因为 $x_{n_j}\to x$, 取 $j$ 足够大, 使得 $n_j\geq N_1$ 且
\[
\rho(x_{n_j},x)<\frac\varepsilon2.
\]
于是对任意 $n\geq N_1$, 有
\[
\rho(x_n,x)
\leq\rho(x_n,x_{n_j})+\rho(x_{n_j},x)
<\varepsilon.
\]
所以 $x_n\to x$。\qedherecustom

\exercise{1.2.3}

设 $F$ 是所有有限支撑的实数列组成的空间, 并令

\[
\rho(x,y)=\sup_{k\ge1}|\xi_k-\eta_k|.
\]

证明 $F$ 不完备, 并指出它的完备化.

\answer

构造
\[
u^{(N)}=\left(1,\frac12,\ldots,\frac1N,0,0,\ldots\right)\in F.
\]
当 $M>N$ 时,
\[
\rho\left(u^{(M)},u^{(N)}\right)
=\sup_{k\geq1}\left|u^{(M)}_k-u^{(N)}_k\right|
=\frac1{N+1}.
\]
因此当 $M,N\to\infty$ 时,
\[
\rho\left(u^{(M)},u^{(N)}\right)
\leq\frac1{\min\{M,N\}+1}\longrightarrow0,
\]
所以 $\{u^{(N)}\}$ 是 $F$ 中的基本列。

若它在 $F$ 中收敛于某个 $u=(u_k)$, 则对每个固定的 $k$, 当 $N\geq k$ 时有 $u^{(N)}_k=1/k$。上确界距离收敛会推出逐坐标收敛, 因而必须有
\[
u_k=\frac1k\qquad(k=1,2,\ldots).
\]
但是 $u$ 有无穷多个非零坐标, 不属于有限支撑空间 $F$, 矛盾。因此 $F$ 不完备。

下面构造它的完备化。令
\[
C_0=\left\{x=(\xi_k):\xi_k\longrightarrow0\right\},
\qquad
\|x\|_\infty=\sup_{k\geq1}|\xi_k|.
\]
每个趋于零的数列有界, 所以 $\|\cdot\|_\infty$ 在 $C_0$ 上有意义。先证 $C_0$ 完备。设 $\{x^{(m)}\}\subset C_0$ 是上确界范数下的基本列, 写成 $x^{(m)}=(\xi^{(m)}_k)$。固定 $k$, 则
\[
|\xi^{(m)}_k-\xi^{(n)}_k|
\leq\|x^{(m)}-x^{(n)}\|_\infty,
\]
所以 $\{\xi^{(m)}_k\}_m$ 是标量基本列。令 $\xi_k$ 为其极限, 并令 $x=(\xi_k)$。

给定 $\varepsilon>0$, 取 $N$ 使 $m,n\geq N$ 时
\[
\|x^{(m)}-x^{(n)}\|_\infty<\frac\varepsilon2.
\]
固定 $m=N$，对每个 $k$ 先在
\[
|\xi^{(N)}_k-\xi^{(n)}_k|<\frac\varepsilon2
\]
中令 $n\to\infty$，得 $|\xi^{(N)}_k-\xi_k|\leq\varepsilon/2$。由于该估计对所有 $k$ 成立，再取上确界，得
\[
\|x^{(N)}-x\|_\infty\leq\frac\varepsilon2.
\]
这也说明 $x$ 有界，因为 $|\xi_k|\leq\|x^{(N)}\|_\infty+\varepsilon/2$。于是对 $m\geq N$,
\[
\|x^{(m)}-x\|_\infty
\leq\|x^{(m)}-x^{(N)}\|_\infty+\|x^{(N)}-x\|_\infty
<\varepsilon.
\]
还需证明 $x\in C_0$。因为 $x^{(N)}\in C_0$, 可取 $K$ 使 $k\geq K$ 时 $|\xi^{(N)}_k|<\varepsilon/2$。结合
\[
|\xi_k-\xi^{(N)}_k|\leq\|x-x^{(N)}\|_\infty\leq\frac\varepsilon2,
\]
得到 $|\xi_k|<\varepsilon$。因此 $\xi_k\to0$, 即 $x\in C_0$。

再证 $F$ 在 $C_0$ 中稠密。对任意 $x=(\xi_k)\in C_0$, 定义截断序列
\[
x^{[N]}=(\xi_1,\ldots,\xi_N,0,0,\ldots)\in F.
\]
则
\[
\|x-x^{[N]}\|_\infty=\sup_{k>N}|\xi_k|\longrightarrow0.
\]
所以 $F$ 等距嵌入 $C_0$ 且在其中稠密, 而 $C_0$ 完备, 因此 $F$ 的完备化就是 $(C_0,\|\cdot\|_\infty)$。\qedherecustom

\exercise{1.2.4}

在 $[0,1]$ 上的多项式空间中定义

\[
\rho(p,q)=\int_0^1|p(x)-q(x)|\,dx.
\]

证明该空间不完备, 并指出它的完备化.

\answer

记 $\mathcal P[0,1]$ 为 $[0,1]$ 上的实系数多项式空间, 并将每个多项式视为 $L^1[0,1]$ 中的等价类。先证 $\rho$ 是距离。非负性、对称性和三角不等式来自绝对值。若
\[
\int_0^1|p(x)-q(x)|\,\mathrm dx=0,
\]
而连续函数 $p-q$ 在某点不为零, 则由连续性可找到一个正长度区间, 使 $|p-q|$ 在该区间上有正下界, 从而积分严格为正, 矛盾。因此 $p=q$。

再证 $\mathcal P[0,1]$ 在 $L^1[0,1]$ 中稠密。给定 $f\in L^1[0,1]$ 和 $\varepsilon>0$, 由连续函数在 $L^1[0,1]$ 中的稠密性, 可取 $g\in C[0,1]$, 使
\[
\|f-g\|_1<\frac\varepsilon2.
\]
再由 Weierstrass 多项式逼近定理, 可取多项式 $p$, 使
\[
\|g-p\|_\infty<\frac\varepsilon2.
\]
由于区间长度为 $1$,
\[
\|g-p\|_1\leq\|g-p\|_\infty<\frac\varepsilon2.
\]
故由三角不等式
\[
\|f-p\|_1\leq\|f-g\|_1+\|g-p\|_1<\varepsilon.
\]
所以 $\mathcal P[0,1]$ 在 $L^1[0,1]$ 中稠密。由于 $L^1[0,1]$ 完备, 根据完备化的唯一性, $\mathcal P[0,1]$ 的完备化等距同构于 $L^1[0,1]$。

最后证多项式空间本身不完备。令
\[
f=\mathbf 1_{[0,1/2]}.
\]
由稠密性, 存在多项式列 $p_n$ 使 $p_n\to f$ 于 $L^1[0,1]$。由
\[
\rho(p_m,p_n)=\|p_m-p_n\|_1
\leq\|p_m-f\|_1+\|f-p_n\|_1\longrightarrow0
\quad(m,n\to\infty),
\]
知 $\{p_n\}$ 是 $\rho$ 下的基本列。若它在 $\mathcal P[0,1]$ 中收敛于多项式 $p$, 则 $L^1$ 极限的唯一性给出 $p=f$ 几乎处处成立。

于是 $p-1$ 在 $(0,1/2)$ 上几乎处处为零, $p$ 在 $(1/2,1)$ 上几乎处处为零。连续函数若在一个区间上几乎处处为零, 则它在该区间上处处为零: 否则某点的非零值由连续性会产生一个正长度区间, 与几乎处处为零矛盾。因此 $p=1$ 在左侧区间恒成立, $p=0$ 在右侧区间恒成立。再由连续性, 同时得到 $p(1/2)=1$ 和 $p(1/2)=0$, 矛盾。故该基本列在多项式空间中没有极限, $\mathcal P[0,1]$ 不完备。\qedherecustom

\exercise{1.2.5}

设 $(\mathcal X,\rho)$ 是完备度量空间. 给定序列 $\{x_n\}$, 若对每个 $\varepsilon>0$ 都存在基本序列 $\{y_n\}$ 使得

\[
\rho(x_n,y_n)<\varepsilon\quad(n=1,2,\ldots),
\]

证明 $\{x_n\}$ 收敛.

\answer

要证明 $\{x_n\}$ 收敛, 先证它是基本列。给定 $\delta>0$, 由题设, 存在基本列 $\{y_n\}$, 使得对所有 $n$ 有
\[
\rho(x_n,y_n)<\frac\delta4.
\]
由于 $\{y_n\}$ 是基本列, 存在 $N$, 使得当 $m,n\geq N$ 时
\[
\rho(y_m,y_n)<\frac\delta2.
\]
因此对任意 $m,n\geq N$, 由三角不等式
\[
\begin{aligned}
\rho(x_m,x_n)
&\leq\rho(x_m,y_m)+\rho(y_m,y_n)+\rho(y_n,x_n)\\
&<\frac\delta4+\frac\delta2+\frac\delta4
=\delta.
\end{aligned}
\]
所以 $\{x_n\}$ 是 $\mathcal X$ 中的基本列。因 $(\mathcal X,\rho)$ 完备, 存在 $x\in\mathcal X$, 使 $x_n\to x$。\qedherecustom

\clearpage
\section*{1.3 紧性}

\exercise{1.3.1}

设 $A$ 是完备度量空间中的子集. 证明 $A$ 序列紧, 当且仅当对任意 $\varepsilon>0$, $A$ 有一个序列紧的 $\varepsilon$-网.

\answer

必要性：若 $A$ 列紧，则由 Hausdorff 定理，$A$ 完全有界。因此对任意 $\varepsilon>0$，$A$ 有一个有限的 $\varepsilon$-网。有限集中的任意序列均有常值子列，所以它自然是列紧的 $\varepsilon$-网。

充分性：给定 $\varepsilon>0$，设 $N$ 是 $A$ 的一个列紧的 $\varepsilon/2$-网。由 Hausdorff 定理，存在 $N$ 的有限 $\varepsilon/2$-网 $N_0$。对任意 $x\in A$，先取 $\xi\in N$，再取 $x_\xi\in N_0$，使
\[
\rho(x,\xi)<\frac\varepsilon2,\qquad
\rho(\xi,x_\xi)<\frac\varepsilon2.
\]
于是
\[
\rho(x,x_\xi)\leq\rho(x,\xi)+\rho(\xi,x_\xi)<\varepsilon.
\]
因此 $N_0$ 是 $A$ 的有限 $\varepsilon$-网，从而 $A$ 完全有界。由于所在度量空间完备，由 Hausdorff 定理，$A$ 列紧。这里子列的极限允许在所在空间中，未要求属于 $A$。\qedherecustom

\exercise{1.3.2}

设 $f$ 是度量空间上紧集的连续函数. 证明 $f$ 有界并且在该紧集上取得上确界和下确界.

\answer

设 $f:M\to\mathbb R$，其中 $M$ 是非空紧集。先证明 $f$ 有上界。反证法，若不然，则对每个正整数 $n$，存在 $x_n\in M$，使 $f(x_n)>n$。由于 $M$ 紧，存在子列 $x_{n_k}\to x_0\in M$。由连续性，
\[
f(x_{n_k})\longrightarrow f(x_0),
\]
但 $f(x_{n_k})>n_k\to+\infty$，矛盾。因此 $f$ 有上界。将同一论证用于 $-f$，得到 $f$ 也有下界，所以 $f$ 有界。

再证能达到上确界。由于 $f(M)$ 是非空有界实数集，由实数的确界性质，存在
\[
\alpha=\sup_{x\in M}f(x)\in\mathbb R.
\]
由上确界定义，对每个正整数 $n$，可取 $x_n\in M$，使
\[
\alpha-\frac1n<f(x_n)\leq\alpha.
\]
又因为 $M$ 紧，存在子列 $x_{n_k}\to x_0\in M$。由连续性与夹逼定理，
\[
f(x_{n_k})\longrightarrow f(x_0),\qquad
f(x_{n_k})\longrightarrow\alpha.
\]
由极限唯一性，$f(x_0)=\alpha$，因此上确界可以达到。

下确界同理：记 $\beta=\inf_{x\in M}f(x)$，取 $y_n\in M$ 使
\[
\beta\leq f(y_n)<\beta+\frac1n.
\]
取收敛子列 $y_{n_k}\to y_0\in M$，由连续性得 $f(y_0)=\beta$。\qedherecustom

\exercise{1.3.3}

证明完全有界集必有界. 再在 $\ell^2$ 中令

\[
E=\{e_k\}_{k=1}^{\infty},\qquad e_k=(0,\ldots,0,1,0,\ldots),
\]

说明有界集不一定全有界.

\answer

若 $A=\varnothing$，有界性显然。设 $A\ne\varnothing$ 且完全有界。令 $\varepsilon=1$，则存在有限的 $1$-网 $N=\{x_1,\ldots,x_n\}$，使
\[
A\subset\bigcup_{i=1}^nB(x_i,1).
\]
对任意 $x\in A$，选取 $i$ 使 $\rho(x,x_i)<1$，则
\[
\rho(x,x_1)\leq\rho(x,x_i)+\rho(x_i,x_1)
<1+\max_{1\leq i\leq n}\rho(x_i,x_1).
\]
右侧是与 $x$ 无关的有限常数，所以 $A$ 有界。

再说明有界集不一定完全有界。在 $\ell^2$ 中，$\|e_k\|_2=1$，所以 $E$ 有界；当 $i\ne j$ 时，
\[
\|e_i-e_j\|_2=\sqrt{1^2+(-1)^2}=\sqrt2.
\]
取 $0<\varepsilon<\sqrt2/2$。若某个半径为 $\varepsilon$ 的球同时包含 $e_i,e_j$，设球心为 $z$，则
\[
\sqrt2=\|e_i-e_j\|_2
\leq\|e_i-z\|_2+\|z-e_j\|_2<2\varepsilon<\sqrt2,
\]
矛盾。因此每个这样的球至多包含 $E$ 中的一个点，有限多个球不能覆盖无限集 $E$。所以 $E$ 不完全有界。\qedherecustom

\exercise{1.3.4}

设 $F_1,F_2$ 是度量空间中的两个紧子集. 证明存在 $x_i\in F_i$ 使得

\[
\rho(F_1,F_2)=\rho(x_1,x_2),
\]

其中

\[
\rho(F_1,F_2)=\inf\{\rho(x,y):x\in F_1,y\in F_2\}.
\]

\answer

设 $F_1,F_2$ 均非空，记 $d=\rho(F_1,F_2)$。由下确界定义，对每个正整数 $n$，可取 $x_n\in F_1$、$y_n\in F_2$，使
\[
d\leq\rho(x_n,y_n)<d+\frac1n.
\]
由 $F_1$ 紧，存在子列 $x_{n_k}\to x_0\in F_1$。再对具有相同指标的序列 $\{y_{n_k}\}$ 使用 $F_2$ 的紧性，取子列
\[
y_{n_{k_j}}\longrightarrow y_0\in F_2.
\]
同时 $\{x_{n_{k_j}}\}$ 仍是收敛列 $\{x_{n_k}\}$ 的子列，所以 $x_{n_{k_j}}\to x_0$。由三角不等式，
\[
\left|\rho(x_{n_{k_j}},y_{n_{k_j}})-\rho(x_0,y_0)\right|
\leq\rho(x_{n_{k_j}},x_0)+\rho(y_{n_{k_j}},y_0)
\longrightarrow0.
\]
另一方面，由最初的不等式及夹逼定理，$\rho(x_{n_{k_j}},y_{n_{k_j}})\to d$。由极限唯一性，$\rho(x_0,y_0)=d$。取 $x_1=x_0\in F_1$、$x_2=y_0\in F_2$ 即得结论。\qedherecustom

\exercise{1.3.5}

设 $M$ 是 $C[a,b]$ 中的有界集, 定义

\[
F(x)=\left\{\int_a^x f(t)\,dt:f\in M\right\}.
\]

证明函数族

\[
F=\left\{F_f:F_f(x)=\int_a^x f(t)\,dt,\ f\in M\right\}
\]

是列紧集.

\answer

使用上确界度量
\[
d(u,v)=\max_{a\leq t\leq b}|u(t)-v(t)|.
\]
若 $M=\varnothing$，结论显然。以下设 $M\ne\varnothing$。先证 $F$ 一致有界。由 $M$ 在 $C[a,b]$ 中有界，可取常数 $m>0$，使
\[
\|f\|_\infty\leq m\qquad(f\in M).
\]
对任意 $f\in M$ 及 $x\in[a,b]$，有
\[
|F_f(x)|=\left|\int_a^xf(t)\,dt\right|
\leq\int_a^x|f(t)|\,dt\leq m(x-a)\leq m(b-a).
\]
所以 $\|F_f\|_\infty\leq m(b-a)$ 对所有 $f\in M$ 成立，即 $F$ 一致有界。

再证等度连续。对任意 $f\in M$ 和 $x,y\in[a,b]$，有
\[
|F_f(x)-F_f(y)|=\left|\int_y^xf(t)\,dt\right|\leq m|x-y|.
\]
给定 $\varepsilon>0$，取 $\delta=\varepsilon/m$，则只要 $|x-y|<\delta$，就对所有 $f\in M$ 同时有 $|F_f(x)-F_f(y)|<\varepsilon$。因此 $F$ 等度连续。

由 Arzel\`a--Ascoli 定理，$F$ 中任意序列均有一个在 $[a,b]$ 上一致收敛的子列，其极限属于 $C[a,b]$。因此 $F$ 列紧。这里不要求极限仍属于 $F$。\qedherecustom

\end{document}
